\documentclass[11pt,a4paper]{article}

\usepackage{kotex}
\usepackage{fontspec}
\setmainfont{Times New Roman}
\setmainhangulfont{AppleMyungjo}
\setsanshangulfont{Apple SD Gothic Neo}

\usepackage[a4paper,top=18mm,bottom=18mm,left=20mm,right=20mm]{geometry}
\usepackage{fancyhdr}
\setlength{\headheight}{24pt}
\usepackage{enumitem}
\usepackage{mdframed}
\usepackage{array}
\usepackage{booktabs}

\setlength{\parindent}{1em}
\linespread{1.18}

\pagestyle{fancy}
\fancyhf{}
\renewcommand{\headrulewidth}{0.6pt}
\fancyhead[L]{\sffamily\bfseries 2026 고1 영어 변형 --- 정답 및 해설 35}
\fancyhead[R]{\sffamily [디지털 미니멀리즘]}
\renewcommand{\footrulewidth}{0pt}
\fancyfoot[C]{\framebox[16mm][c]{\thepage}}

\newcommand{\qno}[1]{\par\vspace{10pt}\noindent{\sffamily\bfseries #1.}\ }
\newmdenv[linewidth=0.6pt,roundcorner=3pt,innertopmargin=7pt,innerbottommargin=7pt,
          innerleftmargin=9pt,innerrightmargin=9pt,skipabove=5pt,skipbelow=5pt]{pbox}

\begin{document}

\begin{center}
{\fontsize{15}{17}\selectfont\sffamily\bfseries [지문 35] 정답 및 해설}
\end{center}
\vspace{6pt}

%% 정답 표
\begin{center}
\begin{tabular}{cccc}
\toprule
\textbf{문항} & \textbf{유형} & \textbf{정답} & \textbf{배점} \\
\midrule
1 & 무관 문장 & ③ & 2점 \\
2 & 빈칸추론 & ④ & 2점 \\
3 & 영어 주제 & ② & 2점 \\
4 & 서술형 의미 & 모범답안 참조 & 2점 \\
\bottomrule
\end{tabular}
\end{center}

\bigskip

%% 문항 1 해설
\qno{1}\ \textbf{[무관 문장 --- 2점] 정답: ③}

\vspace{4pt}
\noindent\textbf{근거}: 이 글은 디지털 미니멀리즘이 스트레스·불안·우울 감소, 수면 개선 등 \textbf{정신 건강}에 미치는 긍정적 효과를 사례 연구를 통해 서술한다.

\vspace{4pt}
\noindent\textbf{③ 분석}: \textit{``Additionally, participants reported relying on social media to get information for their purchasing decisions.''} --- 이 문장은 소셜 미디어를 \textbf{구매 의사결정 정보 수집 수단}으로 사용한다는 내용으로, 글의 중심 주제인 디지털 미니멀리즘과 정신 건강의 관계와 무관하다. 소셜 미디어를 언급하여 주제와 관련 있어 보이지만, 논리적 흐름(스트레스 감소 $\to$ 집중력 향상 $\to$ 구매 정보? $\to$ 번아웃 극복)이 단절된다.

\vspace{6pt}
\noindent\textbf{오답 소거}:
\begin{itemize}[leftmargin=2em,itemsep=1pt,topsep=2pt]
  \item ① Stanford 연구 사례 도입 --- 스트레스·불안·우울 감소 근거 제시. 흐름 일치.
  \item ② 정보·알림 감소 $\to$ 의미 있는 활동 집중 --- 정신 건강 향상 메커니즘 설명. 흐름 일치.
  \item ④ 번아웃·불면증 직장인의 디지털 미니멀리즘 실천 사례 도입 --- 흐름 일치.
  \item ⑤ 수면 질 향상·스트레스 감소·평온함 회복 --- ④의 결과 서술. 흐름 일치.
\end{itemize}

\bigskip

%% 문항 2 해설
\qno{2}\ \textbf{[빈칸추론 --- 2점] 정답: ④ psychological well-being}

\vspace{4pt}
\noindent\textbf{근거 문장}: \textit{``It is evident through various case studies that adopting digital minimalism can lead to improved \underline{\hspace{3.5cm}}.''} --- 이후 전개되는 두 가지 사례(Stanford 연구의 스트레스·불안·우울 감소, 직장인의 수면 개선·스트레스 감소·평온함 회복)가 모두 \textbf{심리적 안녕(psychological well-being)}을 구체적으로 예시한다.

\vspace{4pt}
\noindent\textbf{오답 소거}:
\begin{itemize}[leftmargin=2em,itemsep=1pt,topsep=2pt]
  \item ① productivity at work --- 업무 생산성은 본문에서 직접 다루지 않는다.
  \item ② digital literacy skills --- 디지털 활용 능력은 글의 논지(사용 제한의 효과)와 반대 방향.
  \item ③ consumer decision-making --- ③번 무관 문장의 소재로, 함정 선지. 글 전체 주제와 무관.
  \item ⑤ social media engagement --- 소셜 미디어 참여 증대는 디지털 미니멀리즘의 목표와 정반대.
\end{itemize}

\bigskip

%% 문항 3 해설
\qno{3}\ \textbf{[영어 주제 --- 2점] 정답: ② evidence that reducing digital usage can benefit mental health}

\vspace{4pt}
\noindent\textbf{해설}: 글은 Stanford 연구와 직장인 사례라는 두 가지 실증적 근거를 통해 디지털 사용 감소가 정신 건강에 미치는 긍정적 효과를 논증한다. ②는 이 구조(증거 $\to$ 정신 건강 혜택)를 정확히 포착한다.

\vspace{6pt}
\noindent\textbf{오답 소거}:
\begin{itemize}[leftmargin=2em,itemsep=1pt,topsep=2pt]
  \item ① the growing popularity of digital minimalism among young consumers --- 소비자·인기도는 본문에 없는 내용.
  \item ③ how social media affects purchasing behavior in the digital age --- ③번 무관 문장의 소재를 주제로 오인하게 만드는 함정 선지.
  \item ④ techniques for practicing mindfulness in a connected world --- 마음챙김은 한 사례의 세부 기법이지 글 전체 주제가 아님.
  \item ⑤ the relationship between screen time and academic performance --- 학업 성취는 본문에 언급되지 않음.
\end{itemize}

\vspace{8pt}
\noindent{\sffamily\bfseries 4. 서술형}
\par\textbf{모범답안}: 휴대폰을 언제, 얼마나, 어떤 상황에서 사용할지 스스로 제한을 두는 것이다.
\par\textbf{허용답안}: 휴대폰 사용 시간이나 사용 상황에 규칙을 정하는 것. / 무분별한 휴대폰 사용을 막기 위해 사용 범위를 정하는 것.
\par\textbf{채점 기준}: ``휴대폰 사용''과 ``제한/규칙/범위 설정''의 의미가 모두 드러나면 정답.

\end{document}
